\section{Posterior laws and uniform local estimates}\label{sec:posterior}

An experiment has a uniform input $I\in[4]$ and an observation $Y$.
Its posterior law is the distribution of
$p=\Prob(I\in\cdot\mid Y)$ under the marginal law of $Y$.
Call this law symmetric if it is invariant under every permutation of
the four coordinates. Define
\[
 L=\begin{pmatrix}1&1&-1&-1\\1&-1&1&-1\\1&-1&-1&1\end{pmatrix},
 \qquad x=Lp,\qquad A(p)=|x|^2.
\]
Since $p=\unif+L^{\mathsf T}x/4$, we have
$A(p)=4|p-\unif|^2$. For a posterior law $\mu$, write
\[
 m(\mu)=\E_\mu A,\qquad
 \tau(\mu)=\E_\mu x_1x_2x_3,\qquad h(\mu)=|\tau(\mu)|.
\]
Always $|x_i|\leq1$, $0\leq m\leq3$, and $h\leq m$.
For a symmetric law, coordinate permutations and even sign changes give
\begin{equation}\label{eq:symmetry}
 \E x_i=0,\qquad \E x_ix_j=\frac m3\mathbf1_{\{i=j\}}.
\end{equation}
Every degree-three coordinate monomial other than $x_1x_2x_3$ has
zero expectation. Every monomial of degree at least two has absolute
value at most $|x|^2$: retain two factors and use $|x_i|\leq1$.

\subsection{Channel steps and products}
A channel step of parameter $t\in[0,1]$ replaces a posterior $p$ by
$\unif+t(p-\unif)$. This follows by composing the experiment with the
symmetric channel $P_t$ and applying Bayes' formula under the uniform
prior. It sends
\begin{equation}\label{eq:scaling}
 (x,m,\tau,h)\longmapsto(tx,t^2m,t^3\tau,t^3h).
\end{equation}

For two experiments independent conditional on the same input, denote
the product posterior law by $\mu\star\nu$. To compute it, first draw
$p\sim\mu,q\sim\nu$ independently under their marginal laws. Set
\begin{equation}\label{eq:bayes}
 Z=4\sum_{i=1}^4p_iq_i,\qquad P_i=\frac{4p_iq_i}{Z}\quad(Z>0).
\end{equation}
The joint observation marginal has density $Z$ with respect to the
product of the two marginals, and $\E Z=1$. Thus
$\E_{\mu\star\nu}f=\E_{\mu\otimes\nu}[Zf(P)]$.
Assign any posterior on $Z=0$, where the weighted contribution is zero.
These formulas also show that channel steps and products preserve
symmetry.

\subsection{The assumed moment saving}
Here is the complete input used from \cite{Threshold,Capacity}.

\begin{assumption}\label{ass:input}
There is a class $\calC$ of symmetric posterior laws containing the
revealing law, uniform on the four point masses, and the uninformative
law $\delta_\unif$. It is closed under channel steps with every
$t\in[0,1]$ and under products of experiments independent conditional
on the same uniform input. For $\mu,\nu\in\calC$,
\begin{equation}\label{eq:saving}
 m(\mu\star\nu)\leq m(\mu)+m(\nu)
 -\frac{\E_\mu[A^2]m(\nu)+m(\mu)\E_\nu[A^2]}{1000}.
\end{equation}
\end{assumption}

The class specified in \cite[Proposition 3.2]{Threshold} consists of
symmetric laws with $\E H(p)\geq0$, where, with
$B=x_1x_2x_3$ and $J=x_1^2x_2^2+x_1^2x_3^2+x_2^2x_3^2$,
\begin{equation}\label{eq:cone}
 H=B-\frac{A^2}{10}-\frac{3J}{20}+\frac{3AB}{10}
       -\frac{A^3}{50}+\frac{7AJ}{100}-\frac{51B^2}{100}.
\end{equation}
That proposition asserts the closure properties and
\eqref{eq:saving}; \cite{Capacity} restates them as its companion
input. We use the repository versions fixed in the bibliography and
do not prove Assumption~\ref{ass:input} here.

Every finite subtree experiment in this paper has law in $\calC$.
At a terminal we observe its spin, possibly after an independent
channel step, or observe nothing. Induction toward the root uses
\eqref{eq:scaling} and conditionally independent products. This also
proves membership for every partial product of child experiments.
No closure under limits is assumed or used.

\subsection{Estimates without smoothing}
The following estimates require only symmetry.

\begin{lemma}\label{lem:local}
Let $\mu,\nu$ be symmetric four-state posterior laws, and put
$a=m(\mu)$, $b=m(\nu)$, $h_\mu=h(\mu)$, $h_\nu=h(\nu)$.
Then
\begin{align}
 m(\mu\star\nu)&\geq a+b-7(h_\mu b+a h_\nu),\label{eq:local-m}\\
 h(\mu\star\nu)&\leq h_\mu+h_\nu+460ab.\label{eq:local-h}
\end{align}
The bounds include laws supported on the boundary of the simplex.
\end{lemma}

\begin{proof}
All expectations below are under independent marginal coordinates
$x=Lp$, $y=Lq$. Put $s=x\cdot y$ and, for
$\{i,j,k\}=\{1,2,3\}$, set
\[
 N_i=x_i+y_i+x_jy_k+x_ky_j.
\]
Equation~\eqref{eq:bayes} gives $Z=1+s$, $LP=N/Z$ on $Z>0$.
In particular,
\begin{equation}\label{eq:boundary}
 0\leq Z\leq4,\qquad |N_i|\leq Z,
 \qquad N=0\text{ when }Z=0.
\end{equation}
The last assertion follows also from the unnormalized Bayes numerator.

Write $N=x+y+W$, where $W_i=x_jy_k+x_ky_j$, and put
$\alpha=\tau(\mu)$, $\gamma=\tau(\nu)$.
Symmetry and independence give
\[
 \E|W|^2=\frac23ab,\quad \E s^2=\frac13ab,\quad
 \E[s\,x\cdot W]=2\alpha b,\quad
 \E[s\,y\cdot W]=2a\gamma,\quad
 \E[s|W|^2]=6\alpha\gamma.
\]
For example, $x\cdot W=2(x_1x_2y_3+x_1x_3y_2+x_2x_3y_1)$;
after multiplication by $s$, only the third moment with three distinct
indices survives in $x$. Expanding $|N|^2(1-s)$, terms of degree one
in either vector vanish, and the terms $\E|W|^2$ and $2\E s^2$
cancel. Consequently
\begin{equation}\label{eq:exact-m}
 m(\mu\star\nu)=a+b-4(\alpha b+a\gamma)-6\alpha\gamma
                  +\E\frac{|N|^2s^2}{1+s}.
\end{equation}
Here we used $(1+s)^{-1}=1-s+s^2/(1+s)$. The last term is
nonnegative, and extends by zero at $Z=0$ because
$|N|^2/Z\leq3Z$. Since
$|\alpha\gamma|\leq(h_\mu b+a h_\nu)/2$, this proves
\eqref{eq:local-m}.

For the third moment,
\[
 \tau(\mu\star\nu)=\E\frac{N_1N_2N_3}{(1+s)^2},\qquad
 \frac1{(1+s)^2}-(1-2s)=\frac{s^2(3+2s)}{(1+s)^2}.
\]
The absolute remainder after multiplication by $N_1N_2N_3$ is at
most $36s^2$, by \eqref{eq:boundary}, since
$|N_1N_2N_3|\leq Z^3$ and $|3+2s|\leq9$. Its expectation
is at most $12ab$. All expressions extend by zero at $Z=0$.

In the polynomial $N_1N_2N_3(1-2s)$ the only pure terms are
$x_1x_2x_3$ and $y_1y_2y_3$. Mixed monomials of degree one in
one vector have zero expectation. Every remaining monomial has
absolute expectation at most $ab$, by the monomial bound above and
independence. The sum of the absolute coefficients is at most
$4^3(1+2\cdot3)=448$. Thus the polynomial expectation differs
from $\alpha+\gamma$ by at most $448ab$. Adding the remainder
and taking absolute values proves \eqref{eq:local-h}.
\end{proof}
